\section{账本事件的区域制度深描：western-jin}

\label{sec:ledger-depth-western-jin}

本组叙事以本卷事件账本的可核年序为骨架，将政治节点放回地方资源、制度中介、人口流动与材料类型中。每条来源能够证明的范围不同，故不以朝廷命令、战报数字或后出传记替代基层社会的完整经验。

\subsection{地方资源与制度接口}

\eventanchor{JH-0281-JI-TOMB-BAMBOO}{太康二年（281年）·汲郡魏襄王墓出土竹书，朝廷命荀勖、和峤等整理其书写材料}账本条目\texttt{JH-0281-JI-TOMB-BAMBOO}记载，太康二年（281年），西晋发生“汲郡魏襄王墓出土竹书，朝廷命荀勖、和峤等整理其书写材料”。这一节点可放在盗发王墓使大批竹简进入官方视野，朝廷需要识读古文字并编目这些异于经传的材料的条件下理解；其后续影响被账本概括为竹书成为西晋知识与校书活动的重要资源，也为后世重建先秦纪年留下传本分歧。就地方尺度而言，事件并不只属于朝廷或统帅：征发、道路、仓储、户籍、市场、寺院与家庭关系都可能决定命令如何抵达、战事如何延长，或接收何以在不同地区呈现不一样的速度。

本条列举的核心材料为《晋书·武帝纪》；《晋书·束皙传》；《荀勖传》；汲郡魏襄王墓出土竹书。这些材料较适合钉住事件的发生、官方表述或特定人物、地点，却不自动构成对全境人口、收入、兵数、信仰或动机的调查。账本的置信与材料提示是“发掘与入藏事实可靠；墓主认定、竹书篇章原貌和今本《竹书纪年》与原简关系有争议”；来源限度进一步说明“西晋秘阁整理、汲冢书传本与战国史书重构研究”。因此，不能把条目中的政权名称、行政分类或战报修辞直接转译为固定族群和统一社会意志，也不能将制度宣布写成各地同步实施。

在叙事上，更稳妥的做法是把这一事件视为一条需要地方中介才能运作的链：中央的命令、地方官与精英的选择、运输与劳力的条件、以及普通家庭可能面对的迁徙、赋役或安全风险彼此相连。现有材料能够显示其中若干环节，却保留了大量沉默。承认这种沉默并非放弃解释，而是避免以一条编年记录覆盖不同区域和社群的差异。

\eventanchor{JH-0289-YANG-JUN-REGENT}{太熙元年（289年）·武帝病重时以外戚杨骏受遗诏辅政，围绕太子司马衷的继承安排集中于宫廷}账本条目\texttt{JH-0289-YANG-JUN-REGENT}记载，太熙元年（289年），西晋发生“武帝病重时以外戚杨骏受遗诏辅政，围绕太子司马衷的继承安排集中于宫廷”。这一节点可放在武帝年老而皇太子政治能力备受质疑，杨氏既掌后族又控制中枢官署的条件下理解；其后续影响被账本概括为外戚辅政的脆弱结构为次年权力冲突预设了条件。就地方尺度而言，事件并不只属于朝廷或统帅：征发、道路、仓储、户籍、市场、寺院与家庭关系都可能决定命令如何抵达、战事如何延长，或接收何以在不同地区呈现不一样的速度。

本条列举的核心材料为《晋书·武帝纪》；《晋书·杨骏传》；《资治通鉴·晋纪》。这些材料较适合钉住事件的发生、官方表述或特定人物、地点，却不自动构成对全境人口、收入、兵数、信仰或动机的调查。账本的置信与材料提示是“受遗辅政与杨骏居要职可靠；遗诏措辞、武帝真实意向和其他顾命大臣的作用有争议”；来源限度进一步说明“西晋外戚辅政、遗诏政治与宫城军权研究”。因此，不能把条目中的政权名称、行政分类或战报修辞直接转译为固定族群和统一社会意志，也不能将制度宣布写成各地同步实施。

在叙事上，更稳妥的做法是把这一事件视为一条需要地方中介才能运作的链：中央的命令、地方官与精英的选择、运输与劳力的条件、以及普通家庭可能面对的迁徙、赋役或安全风险彼此相连。现有材料能够显示其中若干环节，却保留了大量沉默。承认这种沉默并非放弃解释，而是避免以一条编年记录覆盖不同区域和社群的差异。

\eventanchor{JH-0290-WU-EMPEROR-DEATH}{太熙元年四月（290年）·司马炎去世，惠帝司马衷即位，杨骏以太傅录尚书事主持朝政}账本条目\texttt{JH-0290-WU-EMPEROR-DEATH}记载，太熙元年四月（290年），西晋发生“司马炎去世，惠帝司马衷即位，杨骏以太傅录尚书事主持朝政”。这一节点可放在统一期的皇权继承未能建立可约束后族与宗王的稳定程序的条件下理解；其后续影响被账本概括为中央权力迅速转入杨氏与皇后贾氏的竞争，太康秩序开始瓦解。就地方尺度而言，事件并不只属于朝廷或统帅：征发、道路、仓储、户籍、市场、寺院与家庭关系都可能决定命令如何抵达、战事如何延长，或接收何以在不同地区呈现不一样的速度。

本条列举的核心材料为《晋书·武帝纪》；《晋书·惠帝纪》；《晋书·杨骏传》。这些材料较适合钉住事件的发生、官方表述或特定人物、地点，却不自动构成对全境人口、收入、兵数、信仰或动机的调查。账本的置信与材料提示是“卒年、继位与杨骏辅政可靠；惠帝初政的实际决策链和宫中诏令形成过程有争议”；来源限度进一步说明“西晋统一后继承、外戚与官僚权力转换研究”。因此，不能把条目中的政权名称、行政分类或战报修辞直接转译为固定族群和统一社会意志，也不能将制度宣布写成各地同步实施。

在叙事上，更稳妥的做法是把这一事件视为一条需要地方中介才能运作的链：中央的命令、地方官与精英的选择、运输与劳力的条件、以及普通家庭可能面对的迁徙、赋役或安全风险彼此相连。现有材料能够显示其中若干环节，却保留了大量沉默。承认这种沉默并非放弃解释，而是避免以一条编年记录覆盖不同区域和社群的差异。

\eventanchor{JH-0291-JIA-OVERTHROWS-YANG}{永平元年（291年）·贾后联络楚王司马玮等发动政变，杨骏被杀，杨氏外戚集团被清除}账本条目\texttt{JH-0291-JIA-OVERTHROWS-YANG}记载，永平元年（291年），西晋发生“贾后联络楚王司马玮等发动政变，杨骏被杀，杨氏外戚集团被清除”。这一节点可放在杨骏排斥异己并限制贾后参与朝政，皇后遂借宗王和宫卫改变权力格局的条件下理解；其后续影响被账本概括为外戚专权并未被制度性解决，宫廷政变成为宗王争夺中枢的先例。就地方尺度而言，事件并不只属于朝廷或统帅：征发、道路、仓储、户籍、市场、寺院与家庭关系都可能决定命令如何抵达、战事如何延长，或接收何以在不同地区呈现不一样的速度。

本条列举的核心材料为《晋书·惠帝纪》；《晋书·贾后传》；《晋书·杨骏传》；《资治通鉴·晋纪》。这些材料较适合钉住事件的发生、官方表述或特定人物、地点，却不自动构成对全境人口、收入、兵数、信仰或动机的调查。账本的置信与材料提示是“政变、杨骏被杀和杨氏失势可靠；贾后、司马玮与禁军将领各自谋划及诏令真伪有争议”；来源限度进一步说明“贾后政变、宫廷武力与西晋外戚政治研究”。因此，不能把条目中的政权名称、行政分类或战报修辞直接转译为固定族群和统一社会意志，也不能将制度宣布写成各地同步实施。

在叙事上，更稳妥的做法是把这一事件视为一条需要地方中介才能运作的链：中央的命令、地方官与精英的选择、运输与劳力的条件、以及普通家庭可能面对的迁徙、赋役或安全风险彼此相连。现有材料能够显示其中若干环节，却保留了大量沉默。承认这种沉默并非放弃解释，而是避免以一条编年记录覆盖不同区域和社群的差异。

\eventanchor{JH-0291-LIANG-WEI-PURGE}{元康元年（291年）·汝南王司马亮与卫瓘辅政后又被贾后借楚王司马玮诛杀，司马玮旋即被处死}账本条目\texttt{JH-0291-LIANG-WEI-PURGE}记载，元康元年（291年），西晋发生“汝南王司马亮与卫瓘辅政后又被贾后借楚王司马玮诛杀，司马玮旋即被处死”。这一节点可放在杨氏倒台后辅政名义落在宗室与重臣手中，但贾后仍控制宫中信息与诏令渠道的条件下理解；其后续影响被账本概括为连续清洗摧毁了稳定辅政的可能，宗王带兵入京的政治惯例进一步固化。就地方尺度而言，事件并不只属于朝廷或统帅：征发、道路、仓储、户籍、市场、寺院与家庭关系都可能决定命令如何抵达、战事如何延长，或接收何以在不同地区呈现不一样的速度。

本条列举的核心材料为《晋书·惠帝纪》；《晋书·汝南文成王亮传》；《晋书·卫瓘传》；《晋书·楚隐王玮传》。这些材料较适合钉住事件的发生、官方表述或特定人物、地点，却不自动构成对全境人口、收入、兵数、信仰或动机的调查。账本的置信与材料提示是“司马亮、卫瓘遇害及司马玮被杀可靠；先后诏命、贾后责任与辅政集团内部关系有争议”；来源限度进一步说明“八王之乱前夜的顾命体制、诏令与清洗研究”。因此，不能把条目中的政权名称、行政分类或战报修辞直接转译为固定族群和统一社会意志，也不能将制度宣布写成各地同步实施。

在叙事上，更稳妥的做法是把这一事件视为一条需要地方中介才能运作的链：中央的命令、地方官与精英的选择、运输与劳力的条件、以及普通家庭可能面对的迁徙、赋役或安全风险彼此相连。现有材料能够显示其中若干环节，却保留了大量沉默。承认这种沉默并非放弃解释，而是避免以一条编年记录覆盖不同区域和社群的差异。

\eventanchor{JH-0296-QI-WANNIAN}{元康六年至九年（296年至299年）·氐族齐万年在关中起兵，西晋多次征讨并长期调动关陇军粮，至299年叛乱被平定}账本条目\texttt{JH-0296-QI-WANNIAN}记载，元康六年至九年（296年至299年），西晋与关陇氐羌集团发生“氐族齐万年在关中起兵，西晋多次征讨并长期调动关陇军粮，至299年叛乱被平定”。这一节点可放在关陇军镇、迁徙人口和土地资源紧张，中央又缺少持久边防与安置方案的条件下理解；其后续影响被账本概括为战争耗损关中并加深朝臣对边民与迁徙的排斥性讨论，削弱西晋西部防线。就地方尺度而言，事件并不只属于朝廷或统帅：征发、道路、仓储、户籍、市场、寺院与家庭关系都可能决定命令如何抵达、战事如何延长，或接收何以在不同地区呈现不一样的速度。

本条列举的核心材料为《晋书·惠帝纪》；《晋书·张方传》；《晋书·江统传》；《资治通鉴·晋纪》。这些材料较适合钉住事件的发生、官方表述或特定人物、地点，却不自动构成对全境人口、收入、兵数、信仰或动机的调查。账本的置信与材料提示是“叛乱延续、主要将领和齐万年败亡大体可靠；参与部落、兵员数字与地方民众立场有争议”；来源限度进一步说明“西晋关中军费、氐羌政治与《徙戎论》语境研究”。因此，不能把条目中的政权名称、行政分类或战报修辞直接转译为固定族群和统一社会意志，也不能将制度宣布写成各地同步实施。

在叙事上，更稳妥的做法是把这一事件视为一条需要地方中介才能运作的链：中央的命令、地方官与精英的选择、运输与劳力的条件、以及普通家庭可能面对的迁徙、赋役或安全风险彼此相连。现有材料能够显示其中若干环节，却保留了大量沉默。承认这种沉默并非放弃解释，而是避免以一条编年记录覆盖不同区域和社群的差异。

\eventanchor{JH-0300-HEIR-APPARENT-KILLED}{永康元年（300年）·贾后以谋反罪废杀太子司马遹，皇位继承与宫廷联盟由此公开破裂}账本条目\texttt{JH-0300-HEIR-APPARENT-KILLED}记载，永康元年（300年），西晋发生“贾后以谋反罪废杀太子司马遹，皇位继承与宫廷联盟由此公开破裂”。这一节点可放在太子与贾后长期不睦，太子外家和朝臣网络也使皇后感到威胁的条件下理解；其后续影响被账本概括为废杀太子给反贾后宗王以动员名义，直接引发司马伦集团的夺权。就地方尺度而言，事件并不只属于朝廷或统帅：征发、道路、仓储、户籍、市场、寺院与家庭关系都可能决定命令如何抵达、战事如何延长，或接收何以在不同地区呈现不一样的速度。

本条列举的核心材料为《晋书·惠帝纪》；《晋书·愍怀太子遹传》；《晋书·贾后传》。这些材料较适合钉住事件的发生、官方表述或特定人物、地点，却不自动构成对全境人口、收入、兵数、信仰或动机的调查。账本的置信与材料提示是“太子被废杀可靠；伪造诏书、太子是否被诱逼书写和具体执行者有争议”；来源限度进一步说明“西晋皇太子废立、文本伪造与宫廷继承危机研究”。因此，不能把条目中的政权名称、行政分类或战报修辞直接转译为固定族群和统一社会意志，也不能将制度宣布写成各地同步实施。

在叙事上，更稳妥的做法是把这一事件视为一条需要地方中介才能运作的链：中央的命令、地方官与精英的选择、运输与劳力的条件、以及普通家庭可能面对的迁徙、赋役或安全风险彼此相连。现有材料能够显示其中若干环节，却保留了大量沉默。承认这种沉默并非放弃解释，而是避免以一条编年记录覆盖不同区域和社群的差异。

\eventanchor{JH-0301-SIMA-LUN-USURPATION}{永宁元年（301年）·赵王司马伦杀贾后并废惠帝自立为帝，短暂改变西晋帝统}账本条目\texttt{JH-0301-SIMA-LUN-USURPATION}记载，永宁元年（301年），西晋发生“赵王司马伦杀贾后并废惠帝自立为帝，短暂改变西晋帝统”。这一节点可放在贾后废杀太子使宫廷失去支持，司马伦及孙秀利用清君侧名义控制禁军的条件下理解；其后续影响被账本概括为宗王称帝打破皇统禁忌，引起齐王司马冏等诸王联合反击。就地方尺度而言，事件并不只属于朝廷或统帅：征发、道路、仓储、户籍、市场、寺院与家庭关系都可能决定命令如何抵达、战事如何延长，或接收何以在不同地区呈现不一样的速度。

本条列举的核心材料为《晋书·惠帝纪》；《晋书·赵王伦传》；《晋书·贾后传》。这些材料较适合钉住事件的发生、官方表述或特定人物、地点，却不自动构成对全境人口、收入、兵数、信仰或动机的调查。账本的置信与材料提示是“司马伦废立与称帝可靠；其诏册的撰写、朝臣拥立程度和贾后死亡细节有争议”；来源限度进一步说明“司马伦篡位、禅让礼仪与宗王政治合法性研究”。因此，不能把条目中的政权名称、行政分类或战报修辞直接转译为固定族群和统一社会意志，也不能将制度宣布写成各地同步实施。

在叙事上，更稳妥的做法是把这一事件视为一条需要地方中介才能运作的链：中央的命令、地方官与精英的选择、运输与劳力的条件、以及普通家庭可能面对的迁徙、赋役或安全风险彼此相连。现有材料能够显示其中若干环节，却保留了大量沉默。承认这种沉默并非放弃解释，而是避免以一条编年记录覆盖不同区域和社群的差异。

\eventanchor{JH-0301-SIMA-JIONG-COUP}{永宁元年（301年）·齐王司马冏等起兵推翻司马伦，复立惠帝并以大司马辅政}账本条目\texttt{JH-0301-SIMA-JIONG-COUP}记载，永宁元年（301年），西晋发生“齐王司马冏等起兵推翻司马伦，复立惠帝并以大司马辅政”。这一节点可放在司马伦僭号触犯多数宗室与官僚的利益，地方军镇得以借复辟集结的条件下理解；其后续影响被账本概括为惠帝虽复位却无自主军权，辅政宗王之间的竞争转为公开军事冲突。就地方尺度而言，事件并不只属于朝廷或统帅：征发、道路、仓储、户籍、市场、寺院与家庭关系都可能决定命令如何抵达、战事如何延长，或接收何以在不同地区呈现不一样的速度。

本条列举的核心材料为《晋书·惠帝纪》；《晋书·齐王冏传》；《晋书·赵王伦传》；《资治通鉴·晋纪》。这些材料较适合钉住事件的发生、官方表述或特定人物、地点，却不自动构成对全境人口、收入、兵数、信仰或动机的调查。账本的置信与材料提示是“司马伦败死、惠帝复位和司马冏辅政可靠；诸王兵力、洛阳城内响应和战后处置有争议”；来源限度进一步说明“八王之乱早期联盟、都城战事与复辟政治研究”。因此，不能把条目中的政权名称、行政分类或战报修辞直接转译为固定族群和统一社会意志，也不能将制度宣布写成各地同步实施。

在叙事上，更稳妥的做法是把这一事件视为一条需要地方中介才能运作的链：中央的命令、地方官与精英的选择、运输与劳力的条件、以及普通家庭可能面对的迁徙、赋役或安全风险彼此相连。现有材料能够显示其中若干环节，却保留了大量沉默。承认这种沉默并非放弃解释，而是避免以一条编年记录覆盖不同区域和社群的差异。

\eventanchor{JH-0302-SIMA-YING-DEFEATS-JIONG}{太安元年（302年）·河间王司马颙与成都王司马颖起兵，长沙王司马乂在洛阳杀齐王司马冏并接掌朝政}账本条目\texttt{JH-0302-SIMA-YING-DEFEATS-JIONG}记载，太安元年（302年），西晋发生“河间王司马颙与成都王司马颖起兵，长沙王司马乂在洛阳杀齐王司马冏并接掌朝政”。这一节点可放在司马冏专权引起宗王不满，关中与邺城的军事资源被引入洛阳政争的条件下理解；其后续影响被账本概括为辅政权在数月间再度易手，中央对地方军队的依赖进一步加深。就地方尺度而言，事件并不只属于朝廷或统帅：征发、道路、仓储、户籍、市场、寺院与家庭关系都可能决定命令如何抵达、战事如何延长，或接收何以在不同地区呈现不一样的速度。

本条列举的核心材料为《晋书·惠帝纪》；《晋书·齐王冏传》；《晋书·长沙厉王乂传》；《晋书·成都王颖传》。这些材料较适合钉住事件的发生、官方表述或特定人物、地点，却不自动构成对全境人口、收入、兵数、信仰或动机的调查。账本的置信与材料提示是“司马冏被杀和司马乂掌政可靠；司马颙、司马颖军队的协调和贾后余党作用有争议”；来源限度进一步说明“八王之乱的联盟重组、洛阳禁军与辅政名义研究”。因此，不能把条目中的政权名称、行政分类或战报修辞直接转译为固定族群和统一社会意志，也不能将制度宣布写成各地同步实施。

在叙事上，更稳妥的做法是把这一事件视为一条需要地方中介才能运作的链：中央的命令、地方官与精英的选择、运输与劳力的条件、以及普通家庭可能面对的迁徙、赋役或安全风险彼此相连。现有材料能够显示其中若干环节，却保留了大量沉默。承认这种沉默并非放弃解释，而是避免以一条编年记录覆盖不同区域和社群的差异。

\eventanchor{JH-0303-SIMA-AI-DEFENDS-LUOYANG}{太安二年（303年）·长沙王司马乂据洛阳抵抗河间王司马颙和成都王司马颖联军，京师陷入长期战斗}账本条目\texttt{JH-0303-SIMA-AI-DEFENDS-LUOYANG}记载，太安二年（303年），西晋发生“长沙王司马乂据洛阳抵抗河间王司马颙和成都王司马颖联军，京师陷入长期战斗”。这一节点可放在司马乂控制惠帝和洛阳官署，司马颙、司马颖则试图以外军重置辅政的条件下理解；其后续影响被账本概括为都城战火消耗财政和人口，西晋无法有效回应边地与地方危机。就地方尺度而言，事件并不只属于朝廷或统帅：征发、道路、仓储、户籍、市场、寺院与家庭关系都可能决定命令如何抵达、战事如何延长，或接收何以在不同地区呈现不一样的速度。

本条列举的核心材料为《晋书·惠帝纪》；《晋书·长沙厉王乂传》；《晋书·河间平王颙传》；《资治通鉴·晋纪》。这些材料较适合钉住事件的发生、官方表述或特定人物、地点，却不自动构成对全境人口、收入、兵数、信仰或动机的调查。账本的置信与材料提示是“司马乂守城与联军进攻可靠；城内粮价、伤亡及惠帝是否亲自指挥有争议”；来源限度进一步说明“八王之乱的洛阳围城、粮食供给与城市破坏研究”。因此，不能把条目中的政权名称、行政分类或战报修辞直接转译为固定族群和统一社会意志，也不能将制度宣布写成各地同步实施。

在叙事上，更稳妥的做法是把这一事件视为一条需要地方中介才能运作的链：中央的命令、地方官与精英的选择、运输与劳力的条件、以及普通家庭可能面对的迁徙、赋役或安全风险彼此相连。现有材料能够显示其中若干环节，却保留了大量沉默。承认这种沉默并非放弃解释，而是避免以一条编年记录覆盖不同区域和社群的差异。

\subsection{道路、边地与人口移动}

\eventanchor{JH-0303-LI-XIONG-TAKES-CHENGDU}{太安二年（303年）·巴氐首领李雄攻取成都，益州的西晋统治被成汉政权取代}账本条目\texttt{JH-0303-LI-XIONG-TAKES-CHENGDU}记载，太安二年（303年），成汉与益州西晋守军发生“巴氐首领李雄攻取成都，益州的西晋统治被成汉政权取代”。这一节点可放在李特集团在蜀地流民与地方矛盾中壮大，西晋益州军政又因内乱难获援助的条件下理解；其后续影响被账本概括为成都成为成汉都城，西晋失去西南重要税源、兵源和交通节点。就地方尺度而言，事件并不只属于朝廷或统帅：征发、道路、仓储、户籍、市场、寺院与家庭关系都可能决定命令如何抵达、战事如何延长，或接收何以在不同地区呈现不一样的速度。

本条列举的核心材料为《晋书·李雄载记》；《华阳国志·李特雄期寿势志》；《资治通鉴·晋纪》。这些材料较适合钉住事件的发生、官方表述或特定人物、地点，却不自动构成对全境人口、收入、兵数、信仰或动机的调查。账本的置信与材料提示是“取成都与李雄建立控制可靠；战斗时日、地方豪族支持和流民构成有争议”；来源限度进一步说明“成汉建国、巴蜀流民与地方军政重组研究”。因此，不能把条目中的政权名称、行政分类或战报修辞直接转译为固定族群和统一社会意志，也不能将制度宣布写成各地同步实施。

在叙事上，更稳妥的做法是把这一事件视为一条需要地方中介才能运作的链：中央的命令、地方官与精英的选择、运输与劳力的条件、以及普通家庭可能面对的迁徙、赋役或安全风险彼此相连。现有材料能够显示其中若干环节，却保留了大量沉默。承认这种沉默并非放弃解释，而是避免以一条编年记录覆盖不同区域和社群的差异。

\eventanchor{JH-0304-SIMA-AI-KILLED}{永安元年（304年）·司马乂被东海王司马越等交给河间王司马颙军并遭杀害，洛阳政权落入外军支配}账本条目\texttt{JH-0304-SIMA-AI-KILLED}记载，永安元年（304年），西晋发生“司马乂被东海王司马越等交给河间王司马颙军并遭杀害，洛阳政权落入外军支配”。这一节点可放在长期守城使洛阳疲敝，司马越等担忧司马乂独揽政权并选择与外军妥协的条件下理解；其后续影响被账本概括为惠帝被迫移置与各王争夺升级，中央威信继续流失。就地方尺度而言，事件并不只属于朝廷或统帅：征发、道路、仓储、户籍、市场、寺院与家庭关系都可能决定命令如何抵达、战事如何延长，或接收何以在不同地区呈现不一样的速度。

本条列举的核心材料为《晋书·惠帝纪》；《晋书·长沙厉王乂传》；《晋书·东海孝献王越传》。这些材料较适合钉住事件的发生、官方表述或特定人物、地点，却不自动构成对全境人口、收入、兵数、信仰或动机的调查。账本的置信与材料提示是“司马乂被害与司马颙军入洛阳可靠；司马越倒戈动机、交付程序和死亡细节有争议”；来源限度进一步说明“八王之乱的背盟、宫廷武力与都城控制研究”。因此，不能把条目中的政权名称、行政分类或战报修辞直接转译为固定族群和统一社会意志，也不能将制度宣布写成各地同步实施。

在叙事上，更稳妥的做法是把这一事件视为一条需要地方中介才能运作的链：中央的命令、地方官与精英的选择、运输与劳力的条件、以及普通家庭可能面对的迁徙、赋役或安全风险彼此相连。现有材料能够显示其中若干环节，却保留了大量沉默。承认这种沉默并非放弃解释，而是避免以一条编年记录覆盖不同区域和社群的差异。

\eventanchor{JH-0304-LIU-YUAN-HAN}{永兴元年（304年）·刘渊在左国城起兵，自称汉王，汇聚匈奴五部及并州反晋力量}账本条目\texttt{JH-0304-LIU-YUAN-HAN}记载，永兴元年（304年），汉赵与并州诸部发生“刘渊在左国城起兵，自称汉王，汇聚匈奴五部及并州反晋力量”。这一节点可放在西晋内战抽空北方防务，刘渊既有部众网络又可借汉室名号争取地方支持的条件下理解；其后续影响被账本概括为北方出现以平阳为核心的反晋政权，西晋的内乱转为多政权战争。就地方尺度而言，事件并不只属于朝廷或统帅：征发、道路、仓储、户籍、市场、寺院与家庭关系都可能决定命令如何抵达、战事如何延长，或接收何以在不同地区呈现不一样的速度。

本条列举的核心材料为《晋书·刘元海载记》；《资治通鉴·晋纪》；《十六国春秋》辑佚。这些材料较适合钉住事件的发生、官方表述或特定人物、地点，却不自动构成对全境人口、收入、兵数、信仰或动机的调查。账本的置信与材料提示是“刘渊称王和政权初建可靠；五部匈奴的动员范围、汉室谱系说法与部众身份有争议”；来源限度进一步说明“刘渊政权、并州军户与匈奴五部政治研究”。因此，不能把条目中的政权名称、行政分类或战报修辞直接转译为固定族群和统一社会意志，也不能将制度宣布写成各地同步实施。

在叙事上，更稳妥的做法是把这一事件视为一条需要地方中介才能运作的链：中央的命令、地方官与精英的选择、运输与劳力的条件、以及普通家庭可能面对的迁徙、赋役或安全风险彼此相连。现有材料能够显示其中若干环节，却保留了大量沉默。承认这种沉默并非放弃解释，而是避免以一条编年记录覆盖不同区域和社群的差异。

\eventanchor{JH-0304-CHEN-MIN-REBELLION}{永兴元年（304年）·陈敏据江东起兵并自置官属，试图利用八王之乱切割东南地区}账本条目\texttt{JH-0304-CHEN-MIN-REBELLION}记载，永兴元年（304年），西晋与江东地方集团发生“陈敏据江东起兵并自置官属，试图利用八王之乱切割东南地区”。这一节点可放在洛阳反复易手使江东守令难获中央支持，陈敏以地方军事和流民为基础扩张的条件下理解；其后续影响被账本概括为其失败显示江东精英尚未完全脱离晋廷，却预示东晋建国的区域条件。就地方尺度而言，事件并不只属于朝廷或统帅：征发、道路、仓储、户籍、市场、寺院与家庭关系都可能决定命令如何抵达、战事如何延长，或接收何以在不同地区呈现不一样的速度。

本条列举的核心材料为《晋书·陈敏传》；《资治通鉴·晋纪》；《建康实录》。这些材料较适合钉住事件的发生、官方表述或特定人物、地点，却不自动构成对全境人口、收入、兵数、信仰或动机的调查。账本的置信与材料提示是“陈敏割据、任官与最终失败可靠；江东士族的实际支持、军队规模和地方秩序有争议”；来源限度进一步说明“西晋末江东割据、士族选择与区域军政研究”。因此，不能把条目中的政权名称、行政分类或战报修辞直接转译为固定族群和统一社会意志，也不能将制度宣布写成各地同步实施。

在叙事上，更稳妥的做法是把这一事件视为一条需要地方中介才能运作的链：中央的命令、地方官与精英的选择、运输与劳力的条件、以及普通家庭可能面对的迁徙、赋役或安全风险彼此相连。现有材料能够显示其中若干环节，却保留了大量沉默。承认这种沉默并非放弃解释，而是避免以一条编年记录覆盖不同区域和社群的差异。

\eventanchor{JH-0305-SIMA-YUE-COALITION}{永兴二年（305年）·东海王司马越以讨河间王司马颙为名起兵，八王之乱扩大为跨州郡的长期战争}账本条目\texttt{JH-0305-SIMA-YUE-COALITION}记载，永兴二年（305年），西晋发生“东海王司马越以讨河间王司马颙为名起兵，八王之乱扩大为跨州郡的长期战争”。这一节点可放在司马颙挟持惠帝和控制洛阳引起司马越及山东、河北军事集团反弹的条件下理解；其后续影响被账本概括为更多州郡军队被卷入内战，北方防线与漕运系统持续失修。就地方尺度而言，事件并不只属于朝廷或统帅：征发、道路、仓储、户籍、市场、寺院与家庭关系都可能决定命令如何抵达、战事如何延长，或接收何以在不同地区呈现不一样的速度。

本条列举的核心材料为《晋书·惠帝纪》；《晋书·东海孝献王越传》；《晋书·河间平王颙传》。这些材料较适合钉住事件的发生、官方表述或特定人物、地点，却不自动构成对全境人口、收入、兵数、信仰或动机的调查。账本的置信与材料提示是“司马越起兵与主要对手可靠；各州响应次序、兵额和地方被迫参战程度有争议”；来源限度进一步说明“司马越集团、州郡动员与八王之乱后期研究”。因此，不能把条目中的政权名称、行政分类或战报修辞直接转译为固定族群和统一社会意志，也不能将制度宣布写成各地同步实施。

在叙事上，更稳妥的做法是把这一事件视为一条需要地方中介才能运作的链：中央的命令、地方官与精英的选择、运输与劳力的条件、以及普通家庭可能面对的迁徙、赋役或安全风险彼此相连。现有材料能够显示其中若干环节，却保留了大量沉默。承认这种沉默并非放弃解释，而是避免以一条编年记录覆盖不同区域和社群的差异。

\eventanchor{JH-0306-HUI-EMPEROR-DEATH}{光熙元年（306年）·惠帝在洛阳去世，司马炽即位为怀帝，司马越继续掌握主要军政权}账本条目\texttt{JH-0306-HUI-EMPEROR-DEATH}记载，光熙元年（306年），西晋发生“惠帝在洛阳去世，司马炽即位为怀帝，司马越继续掌握主要军政权”。这一节点可放在八王之乱已使惠帝长期受制于宗王，继承只能在司马越军事优势下完成的条件下理解；其后续影响被账本概括为怀帝虽即位却无法整合诸王和州镇，西晋面对汉赵时缺乏统一指挥。就地方尺度而言，事件并不只属于朝廷或统帅：征发、道路、仓储、户籍、市场、寺院与家庭关系都可能决定命令如何抵达、战事如何延长，或接收何以在不同地区呈现不一样的速度。

本条列举的核心材料为《晋书·惠帝纪》；《晋书·怀帝纪》；《晋书·东海孝献王越传》。这些材料较适合钉住事件的发生、官方表述或特定人物、地点，却不自动构成对全境人口、收入、兵数、信仰或动机的调查。账本的置信与材料提示是“惠帝卒与怀帝继位可靠；惠帝死因、司马越的责任和新帝实际权力有争议”；来源限度进一步说明“惠帝死亡、司马越辅政与西晋皇权空洞化研究”。因此，不能把条目中的政权名称、行政分类或战报修辞直接转译为固定族群和统一社会意志，也不能将制度宣布写成各地同步实施。

在叙事上，更稳妥的做法是把这一事件视为一条需要地方中介才能运作的链：中央的命令、地方官与精英的选择、运输与劳力的条件、以及普通家庭可能面对的迁徙、赋役或安全风险彼此相连。现有材料能够显示其中若干环节，却保留了大量沉默。承认这种沉默并非放弃解释，而是避免以一条编年记录覆盖不同区域和社群的差异。

\eventanchor{JH-0306-LI-XIONG-EMPEROR}{晏平元年（306年）·李雄在成都称帝，国号成，使巴蜀割据由起事集团转为帝国体制}账本条目\texttt{JH-0306-LI-XIONG-EMPEROR}记载，晏平元年（306年），成汉发生“李雄在成都称帝，国号成，使巴蜀割据由起事集团转为帝国体制”。这一节点可放在李雄已控制成都及周边，需以官号和年号安置部众并争取地方官僚的条件下理解；其后续影响被账本概括为成汉成为西晋覆亡后持续最久的西南政权，为蜀地提供不同于北方战乱的政治中心。就地方尺度而言，事件并不只属于朝廷或统帅：征发、道路、仓储、户籍、市场、寺院与家庭关系都可能决定命令如何抵达、战事如何延长，或接收何以在不同地区呈现不一样的速度。

本条列举的核心材料为《晋书·李雄载记》；《华阳国志·李特雄期寿势志》；《资治通鉴·晋纪》。这些材料较适合钉住事件的发生、官方表述或特定人物、地点，却不自动构成对全境人口、收入、兵数、信仰或动机的调查。账本的置信与材料提示是“称帝与年号大体可靠；国号使用、礼制规模和成汉对蜀地的实际控制范围有争议”；来源限度进一步说明“成汉帝制、巴蜀地方精英与流民政权合法化研究”。因此，不能把条目中的政权名称、行政分类或战报修辞直接转译为固定族群和统一社会意志，也不能将制度宣布写成各地同步实施。

在叙事上，更稳妥的做法是把这一事件视为一条需要地方中介才能运作的链：中央的命令、地方官与精英的选择、运输与劳力的条件、以及普通家庭可能面对的迁徙、赋役或安全风险彼此相连。现有材料能够显示其中若干环节，却保留了大量沉默。承认这种沉默并非放弃解释，而是避免以一条编年记录覆盖不同区域和社群的差异。

\eventanchor{JH-0308-LIU-YUAN-EMPEROR}{永凤元年（308年）·刘渊称帝，正式以汉为国号并以平阳为政治中心}账本条目\texttt{JH-0308-LIU-YUAN-EMPEROR}记载，永凤元年（308年），汉赵发生“刘渊称帝，正式以汉为国号并以平阳为政治中心”。这一节点可放在刘渊已获得并州与河东多支武装支持，需要将联盟转化为可征发的政权的条件下理解；其后续影响被账本概括为汉赵的帝号与官署使其更易吸纳反晋力量，并直接挑战西晋正统。就地方尺度而言，事件并不只属于朝廷或统帅：征发、道路、仓储、户籍、市场、寺院与家庭关系都可能决定命令如何抵达、战事如何延长，或接收何以在不同地区呈现不一样的速度。

本条列举的核心材料为《晋书·刘元海载记》；《资治通鉴·晋纪》；《十六国春秋》辑佚。这些材料较适合钉住事件的发生、官方表述或特定人物、地点，却不自动构成对全境人口、收入、兵数、信仰或动机的调查。账本的置信与材料提示是“称帝、年号与平阳中心地位可靠；官制仿汉的实际程度及汉室继承叙事有争议”；来源限度进一步说明“汉赵帝号、政治记忆与十六国建国合法性研究”。因此，不能把条目中的政权名称、行政分类或战报修辞直接转译为固定族群和统一社会意志，也不能将制度宣布写成各地同步实施。

在叙事上，更稳妥的做法是把这一事件视为一条需要地方中介才能运作的链：中央的命令、地方官与精英的选择、运输与劳力的条件、以及普通家庭可能面对的迁徙、赋役或安全风险彼此相连。现有材料能够显示其中若干环节，却保留了大量沉默。承认这种沉默并非放弃解释，而是避免以一条编年记录覆盖不同区域和社群的差异。

\eventanchor{JH-0309-HAN-ZHAO-SIEGE}{永嘉三年（309年）·刘聪、王弥等围攻洛阳，西晋京师依靠临时征兵和城防勉强守住}账本条目\texttt{JH-0309-HAN-ZHAO-SIEGE}记载，永嘉三年（309年），汉赵与西晋发生“刘聪、王弥等围攻洛阳，西晋京师依靠临时征兵和城防勉强守住”。这一节点可放在司马越与怀帝集团相互猜忌，北方州镇难以协调援军，汉赵得以直逼京师的条件下理解；其后续影响被账本概括为围城暴露洛阳补给和指挥体系的崩坏，为311年城陷创造条件。就地方尺度而言，事件并不只属于朝廷或统帅：征发、道路、仓储、户籍、市场、寺院与家庭关系都可能决定命令如何抵达、战事如何延长，或接收何以在不同地区呈现不一样的速度。

本条列举的核心材料为《晋书·怀帝纪》；《晋书·刘聪载记》；《晋书·苟晞传》；《资治通鉴·晋纪》。这些材料较适合钉住事件的发生、官方表述或特定人物、地点，却不自动构成对全境人口、收入、兵数、信仰或动机的调查。账本的置信与材料提示是“围攻与洛阳危机可靠；参战兵数、援军行动和围城损失有争议”；来源限度进一步说明“永嘉前洛阳围城、军粮与西晋末年城市防御研究”。因此，不能把条目中的政权名称、行政分类或战报修辞直接转译为固定族群和统一社会意志，也不能将制度宣布写成各地同步实施。

在叙事上，更稳妥的做法是把这一事件视为一条需要地方中介才能运作的链：中央的命令、地方官与精英的选择、运输与劳力的条件、以及普通家庭可能面对的迁徙、赋役或安全风险彼此相连。现有材料能够显示其中若干环节，却保留了大量沉默。承认这种沉默并非放弃解释，而是避免以一条编年记录覆盖不同区域和社群的差异。

\eventanchor{JH-0310-LIU-CONG-ACCESSION}{河瑞二年（310年）·刘渊去世后刘聪经宫廷斗争继位，汉赵继续对西晋发动攻势}账本条目\texttt{JH-0310-LIU-CONG-ACCESSION}记载，河瑞二年（310年），汉赵发生“刘渊去世后刘聪经宫廷斗争继位，汉赵继续对西晋发动攻势”。这一节点可放在刘渊的继承安排未消除诸子竞争，掌军的刘聪需迅速确立宫廷优势的条件下理解；其后续影响被账本概括为刘聪保住政权后将资源集中于洛阳方向，西晋失去喘息机会。就地方尺度而言，事件并不只属于朝廷或统帅：征发、道路、仓储、户籍、市场、寺院与家庭关系都可能决定命令如何抵达、战事如何延长，或接收何以在不同地区呈现不一样的速度。

本条列举的核心材料为《晋书·刘元海载记》；《晋书·刘聪载记》；《资治通鉴·晋纪》。这些材料较适合钉住事件的发生、官方表述或特定人物、地点，却不自动构成对全境人口、收入、兵数、信仰或动机的调查。账本的置信与材料提示是“刘渊卒、刘聪即位与皇位竞争可靠；刘和之变的细节、宗室支持和军队控制关系有争议”；来源限度进一步说明“汉赵继承危机、宗室军权与对晋战争连续性研究”。因此，不能把条目中的政权名称、行政分类或战报修辞直接转译为固定族群和统一社会意志，也不能将制度宣布写成各地同步实施。

在叙事上，更稳妥的做法是把这一事件视为一条需要地方中介才能运作的链：中央的命令、地方官与精英的选择、运输与劳力的条件、以及普通家庭可能面对的迁徙、赋役或安全风险彼此相连。现有材料能够显示其中若干环节，却保留了大量沉默。承认这种沉默并非放弃解释，而是避免以一条编年记录覆盖不同区域和社群的差异。

\eventanchor{JH-0310-SIMA-YUE-DEATH}{永嘉四年（310年）·司马越在东出避乱途中去世，其部众和辎重随即失去统一指挥}账本条目\texttt{JH-0310-SIMA-YUE-DEATH}记载，永嘉四年（310年），西晋发生“司马越在东出避乱途中去世，其部众和辎重随即失去统一指挥”。这一节点可放在司马越长期专政却未能击退汉赵，与怀帝冲突后带兵离开洛阳的条件下理解；其后续影响被账本概括为西晋最大的机动军集团瓦解，洛阳的外援和粮道更加空虚。就地方尺度而言，事件并不只属于朝廷或统帅：征发、道路、仓储、户籍、市场、寺院与家庭关系都可能决定命令如何抵达、战事如何延长，或接收何以在不同地区呈现不一样的速度。

本条列举的核心材料为《晋书·怀帝纪》；《晋书·东海孝献王越传》；《晋书·王衍传》。这些材料较适合钉住事件的发生、官方表述或特定人物、地点，却不自动构成对全境人口、收入、兵数、信仰或动机的调查。账本的置信与材料提示是“司马越卒与部众离散可靠；死因、军队数目和其与怀帝最后关系有争议”；来源限度进一步说明“司马越集团瓦解、军队流动与西晋末年权力结构研究”。因此，不能把条目中的政权名称、行政分类或战报修辞直接转译为固定族群和统一社会意志，也不能将制度宣布写成各地同步实施。

在叙事上，更稳妥的做法是把这一事件视为一条需要地方中介才能运作的链：中央的命令、地方官与精英的选择、运输与劳力的条件、以及普通家庭可能面对的迁徙、赋役或安全风险彼此相连。现有材料能够显示其中若干环节，却保留了大量沉默。承认这种沉默并非放弃解释，而是避免以一条编年记录覆盖不同区域和社群的差异。

\subsection{社会关系与行政分类}

\eventanchor{JH-0311-YONGJIA-FALL}{永嘉五年六月（311年）·汉赵军攻陷洛阳，大量官民死散，西晋都城体系遭到毁灭性破坏}账本条目\texttt{JH-0311-YONGJIA-FALL}记载，永嘉五年六月（311年），汉赵与西晋发生“汉赵军攻陷洛阳，大量官民死散，西晋都城体系遭到毁灭性破坏”。这一节点可放在司马越死后洛阳孤立，汉赵军利用西晋内斗与北方州镇失序发动总攻的条件下理解；其后续影响被账本概括为西晋失去首都和官僚中心，北方士民南迁与地方割据显著加速。就地方尺度而言，事件并不只属于朝廷或统帅：征发、道路、仓储、户籍、市场、寺院与家庭关系都可能决定命令如何抵达、战事如何延长，或接收何以在不同地区呈现不一样的速度。

本条列举的核心材料为《晋书·怀帝纪》；《晋书·刘聪载记》；《晋书·王衍传》；《资治通鉴·晋纪》。这些材料较适合钉住事件的发生、官方表述或特定人物、地点，却不自动构成对全境人口、收入、兵数、信仰或动机的调查。账本的置信与材料提示是“洛阳陷落与大规模死散可靠；杀掠数字、宫城毁坏范围和个别殉难叙事有争议”；来源限度进一步说明“永嘉之乱、洛阳城市破坏与北方人口流动研究”。因此，不能把条目中的政权名称、行政分类或战报修辞直接转译为固定族群和统一社会意志，也不能将制度宣布写成各地同步实施。

在叙事上，更稳妥的做法是把这一事件视为一条需要地方中介才能运作的链：中央的命令、地方官与精英的选择、运输与劳力的条件、以及普通家庭可能面对的迁徙、赋役或安全风险彼此相连。现有材料能够显示其中若干环节，却保留了大量沉默。承认这种沉默并非放弃解释，而是避免以一条编年记录覆盖不同区域和社群的差异。

\eventanchor{JH-0311-HUAI-EMPEROR-CAPTURED}{永嘉五年（311年）·怀帝在洛阳陷落时被汉赵俘获并押往平阳，西晋皇统陷入屈辱性危机}账本条目\texttt{JH-0311-HUAI-EMPEROR-CAPTURED}记载，永嘉五年（311年），汉赵与西晋发生“怀帝在洛阳陷落时被汉赵俘获并押往平阳，西晋皇统陷入屈辱性危机”。这一节点可放在洛阳守军崩溃后怀帝无法获得地方军队接应，汉赵意在利用其皇帝身份的条件下理解；其后续影响被账本概括为关中与江东的晋室力量必须另立政治中心，皇统分裂成为现实。就地方尺度而言，事件并不只属于朝廷或统帅：征发、道路、仓储、户籍、市场、寺院与家庭关系都可能决定命令如何抵达、战事如何延长，或接收何以在不同地区呈现不一样的速度。

本条列举的核心材料为《晋书·怀帝纪》；《晋书·刘聪载记》；《资治通鉴·晋纪》。这些材料较适合钉住事件的发生、官方表述或特定人物、地点，却不自动构成对全境人口、收入、兵数、信仰或动机的调查。账本的置信与材料提示是“被俘与押往平阳可靠；出宫路线、侍从人数和俘后礼仪细节有争议”；来源限度进一步说明“被俘皇帝、战败仪式与晋末合法性危机研究”。因此，不能把条目中的政权名称、行政分类或战报修辞直接转译为固定族群和统一社会意志，也不能将制度宣布写成各地同步实施。

在叙事上，更稳妥的做法是把这一事件视为一条需要地方中介才能运作的链：中央的命令、地方官与精英的选择、运输与劳力的条件、以及普通家庭可能面对的迁徙、赋役或安全风险彼此相连。现有材料能够显示其中若干环节，却保留了大量沉默。承认这种沉默并非放弃解释，而是避免以一条编年记录覆盖不同区域和社群的差异。

\eventanchor{JH-0313-HUAI-EMPEROR-KILLED}{建兴元年（313年）·汉赵杀害被俘的怀帝，长安的司马邺被拥立为愍帝}账本条目\texttt{JH-0313-HUAI-EMPEROR-KILLED}记载，建兴元年（313年），汉赵与西晋发生“汉赵杀害被俘的怀帝，长安的司马邺被拥立为愍帝”。这一节点可放在怀帝长期被俘使晋室名义受损，关中官僚和宗室急需保留晋朝继承线的条件下理解；其后续影响被账本概括为愍帝以长安为中心延续西晋残余政权，但其资源远不足以恢复北方。就地方尺度而言，事件并不只属于朝廷或统帅：征发、道路、仓储、户籍、市场、寺院与家庭关系都可能决定命令如何抵达、战事如何延长，或接收何以在不同地区呈现不一样的速度。

本条列举的核心材料为《晋书·怀帝纪》；《晋书·愍帝纪》；《晋书·刘聪载记》。这些材料较适合钉住事件的发生、官方表述或特定人物、地点，却不自动构成对全境人口、收入、兵数、信仰或动机的调查。账本的置信与材料提示是“怀帝遇害与愍帝即位可靠；杀害场景、刘聪动机和长安拥立程序有争议”；来源限度进一步说明“永嘉后皇统重建、长安政权与俘虏政治研究”。因此，不能把条目中的政权名称、行政分类或战报修辞直接转译为固定族群和统一社会意志，也不能将制度宣布写成各地同步实施。

在叙事上，更稳妥的做法是把这一事件视为一条需要地方中介才能运作的链：中央的命令、地方官与精英的选择、运输与劳力的条件、以及普通家庭可能面对的迁徙、赋役或安全风险彼此相连。现有材料能够显示其中若干环节，却保留了大量沉默。承认这种沉默并非放弃解释，而是避免以一条编年记录覆盖不同区域和社群的差异。

\eventanchor{JH-0316-CHANGAN-FALL}{建兴四年十一月（316年）·刘曜攻陷长安，愍帝出降，西晋在北方的最后朝廷覆灭}账本条目\texttt{JH-0316-CHANGAN-FALL}记载，建兴四年十一月（316年），汉赵与西晋发生“刘曜攻陷长安，愍帝出降，西晋在北方的最后朝廷覆灭”。这一节点可放在愍帝政权缺乏粮饷和外援，刘曜控制关中通道并持续围攻的条件下理解；其后续影响被账本概括为晋室政治中心转向建业，北方进入多政权并立与长期战争阶段。就地方尺度而言，事件并不只属于朝廷或统帅：征发、道路、仓储、户籍、市场、寺院与家庭关系都可能决定命令如何抵达、战事如何延长，或接收何以在不同地区呈现不一样的速度。

本条列举的核心材料为《晋书·愍帝纪》；《晋书·刘曜载记》；《资治通鉴·晋纪》。这些材料较适合钉住事件的发生、官方表述或特定人物、地点，却不自动构成对全境人口、收入、兵数、信仰或动机的调查。账本的置信与材料提示是“长安陷落、愍帝降服和西晋终结可靠；围城饥荒、投降条件及关中地方反应有争议”；来源限度进一步说明“西晋灭亡、长安围城与关中地方社会研究”。因此，不能把条目中的政权名称、行政分类或战报修辞直接转译为固定族群和统一社会意志，也不能将制度宣布写成各地同步实施。

在叙事上，更稳妥的做法是把这一事件视为一条需要地方中介才能运作的链：中央的命令、地方官与精英的选择、运输与劳力的条件、以及普通家庭可能面对的迁徙、赋役或安全风险彼此相连。现有材料能够显示其中若干环节，却保留了大量沉默。承认这种沉默并非放弃解释，而是避免以一条编年记录覆盖不同区域和社群的差异。

\eventanchor{JH-LEDGER-001-1}{太康二年（281年）·“汲郡魏襄王墓出土竹书，朝廷命荀勖、和峤等整理其书写材料”的前期动员与资源准备}账本条目\texttt{JH-LEDGER-001-1}记载，太康二年（281年），西晋发生““汲郡魏襄王墓出土竹书，朝廷命荀勖、和峤等整理其书写材料”的前期动员与资源准备”。这一节点可放在当时的权力、资源与信息约束推动相关过程展开的条件下理解；其后续影响被账本概括为这一过程为“汲郡魏襄王墓出土竹书，朝廷命荀勖、和峤等整理其书写材料”提供可观察的条件。就地方尺度而言，事件并不只属于朝廷或统帅：征发、道路、仓储、户籍、市场、寺院与家庭关系都可能决定命令如何抵达、战事如何延长，或接收何以在不同地区呈现不一样的速度。

本条列举的核心材料为《晋书·武帝纪》；《晋书·束皙传》；《荀勖传》；汲郡魏襄王墓出土竹书。这些材料较适合钉住事件的发生、官方表述或特定人物、地点，却不自动构成对全境人口、收入、兵数、信仰或动机的调查。账本的置信与材料提示是“核心事项已有既有账本来源；本分解节点的参与者、执行范围和时间先后仍须逐案核验”；来源限度进一步说明“西晋秘阁整理、汲冢书传本与战国史书重构研究”。因此，不能把条目中的政权名称、行政分类或战报修辞直接转译为固定族群和统一社会意志，也不能将制度宣布写成各地同步实施。

在叙事上，更稳妥的做法是把这一事件视为一条需要地方中介才能运作的链：中央的命令、地方官与精英的选择、运输与劳力的条件、以及普通家庭可能面对的迁徙、赋役或安全风险彼此相连。现有材料能够显示其中若干环节，却保留了大量沉默。承认这种沉默并非放弃解释，而是避免以一条编年记录覆盖不同区域和社群的差异。

\eventanchor{JH-LEDGER-001-2}{太康二年（281年）·“汲郡魏襄王墓出土竹书，朝廷命荀勖、和峤等整理其书写材料”的命令、联盟或地方响应}账本条目\texttt{JH-LEDGER-001-2}记载，太康二年（281年），西晋发生““汲郡魏襄王墓出土竹书，朝廷命荀勖、和峤等整理其书写材料”的命令、联盟或地方响应”。这一节点可放在当时的权力、资源与信息约束推动相关过程展开的条件下理解；其后续影响被账本概括为这一过程为“汲郡魏襄王墓出土竹书，朝廷命荀勖、和峤等整理其书写材料”提供可观察的条件。就地方尺度而言，事件并不只属于朝廷或统帅：征发、道路、仓储、户籍、市场、寺院与家庭关系都可能决定命令如何抵达、战事如何延长，或接收何以在不同地区呈现不一样的速度。

本条列举的核心材料为《晋书·武帝纪》；《晋书·束皙传》；《荀勖传》；汲郡魏襄王墓出土竹书。这些材料较适合钉住事件的发生、官方表述或特定人物、地点，却不自动构成对全境人口、收入、兵数、信仰或动机的调查。账本的置信与材料提示是“核心事项已有既有账本来源；本分解节点的参与者、执行范围和时间先后仍须逐案核验”；来源限度进一步说明“西晋秘阁整理、汲冢书传本与战国史书重构研究”。因此，不能把条目中的政权名称、行政分类或战报修辞直接转译为固定族群和统一社会意志，也不能将制度宣布写成各地同步实施。

在叙事上，更稳妥的做法是把这一事件视为一条需要地方中介才能运作的链：中央的命令、地方官与精英的选择、运输与劳力的条件、以及普通家庭可能面对的迁徙、赋役或安全风险彼此相连。现有材料能够显示其中若干环节，却保留了大量沉默。承认这种沉默并非放弃解释，而是避免以一条编年记录覆盖不同区域和社群的差异。

\eventanchor{JH-LEDGER-001-3}{太康二年（281年）·汲郡魏襄王墓出土竹书，朝廷命荀勖、和峤等整理其书写材料}账本条目\texttt{JH-LEDGER-001-3}记载，太康二年（281年），西晋发生“汲郡魏襄王墓出土竹书，朝廷命荀勖、和峤等整理其书写材料”。这一节点可放在“汲郡魏襄王墓出土竹书，朝廷命荀勖、和峤等整理其书写材料”发生的政治与区域条件共同作用的条件下理解；其后续影响被账本概括为该节点改变后续资源、联盟或制度处置的条件。就地方尺度而言，事件并不只属于朝廷或统帅：征发、道路、仓储、户籍、市场、寺院与家庭关系都可能决定命令如何抵达、战事如何延长，或接收何以在不同地区呈现不一样的速度。

本条列举的核心材料为《晋书·武帝纪》；《晋书·束皙传》；《荀勖传》；汲郡魏襄王墓出土竹书。这些材料较适合钉住事件的发生、官方表述或特定人物、地点，却不自动构成对全境人口、收入、兵数、信仰或动机的调查。账本的置信与材料提示是“核心事项已有既有账本来源；本分解节点的参与者、执行范围和时间先后仍须逐案核验”；来源限度进一步说明“西晋秘阁整理、汲冢书传本与战国史书重构研究”。因此，不能把条目中的政权名称、行政分类或战报修辞直接转译为固定族群和统一社会意志，也不能将制度宣布写成各地同步实施。

在叙事上，更稳妥的做法是把这一事件视为一条需要地方中介才能运作的链：中央的命令、地方官与精英的选择、运输与劳力的条件、以及普通家庭可能面对的迁徙、赋役或安全风险彼此相连。现有材料能够显示其中若干环节，却保留了大量沉默。承认这种沉默并非放弃解释，而是避免以一条编年记录覆盖不同区域和社群的差异。

\eventanchor{JH-LEDGER-001-4}{太康二年（281年）·“汲郡魏襄王墓出土竹书，朝廷命荀勖、和峤等整理其书写材料”的执行中的空间与人员处置}账本条目\texttt{JH-LEDGER-001-4}记载，太康二年（281年），西晋发生““汲郡魏襄王墓出土竹书，朝廷命荀勖、和峤等整理其书写材料”的执行中的空间与人员处置”。这一节点可放在“汲郡魏襄王墓出土竹书，朝廷命荀勖、和峤等整理其书写材料”发生的政治与区域条件共同作用的条件下理解；其后续影响被账本概括为该节点改变后续资源、联盟或制度处置的条件。就地方尺度而言，事件并不只属于朝廷或统帅：征发、道路、仓储、户籍、市场、寺院与家庭关系都可能决定命令如何抵达、战事如何延长，或接收何以在不同地区呈现不一样的速度。

本条列举的核心材料为《晋书·武帝纪》；《晋书·束皙传》；《荀勖传》；汲郡魏襄王墓出土竹书。这些材料较适合钉住事件的发生、官方表述或特定人物、地点，却不自动构成对全境人口、收入、兵数、信仰或动机的调查。账本的置信与材料提示是“核心事项已有既有账本来源；本分解节点的参与者、执行范围和时间先后仍须逐案核验”；来源限度进一步说明“西晋秘阁整理、汲冢书传本与战国史书重构研究”。因此，不能把条目中的政权名称、行政分类或战报修辞直接转译为固定族群和统一社会意志，也不能将制度宣布写成各地同步实施。

在叙事上，更稳妥的做法是把这一事件视为一条需要地方中介才能运作的链：中央的命令、地方官与精英的选择、运输与劳力的条件、以及普通家庭可能面对的迁徙、赋役或安全风险彼此相连。现有材料能够显示其中若干环节，却保留了大量沉默。承认这种沉默并非放弃解释，而是避免以一条编年记录覆盖不同区域和社群的差异。

\eventanchor{JH-LEDGER-001-5}{太康二年（281年）·“汲郡魏襄王墓出土竹书，朝廷命荀勖、和峤等整理其书写材料”的后续制度或军政调整}账本条目\texttt{JH-LEDGER-001-5}记载，太康二年（281年），西晋发生““汲郡魏襄王墓出土竹书，朝廷命荀勖、和峤等整理其书写材料”的后续制度或军政调整”。这一节点可放在核心节点发生后仍须处理其遗留问题的条件下理解；其后续影响被账本概括为后续影响须在不同地区和材料类型中分别核验。就地方尺度而言，事件并不只属于朝廷或统帅：征发、道路、仓储、户籍、市场、寺院与家庭关系都可能决定命令如何抵达、战事如何延长，或接收何以在不同地区呈现不一样的速度。

本条列举的核心材料为《晋书·武帝纪》；《晋书·束皙传》；《荀勖传》；汲郡魏襄王墓出土竹书。这些材料较适合钉住事件的发生、官方表述或特定人物、地点，却不自动构成对全境人口、收入、兵数、信仰或动机的调查。账本的置信与材料提示是“核心事项已有既有账本来源；本分解节点的参与者、执行范围和时间先后仍须逐案核验”；来源限度进一步说明“西晋秘阁整理、汲冢书传本与战国史书重构研究”。因此，不能把条目中的政权名称、行政分类或战报修辞直接转译为固定族群和统一社会意志，也不能将制度宣布写成各地同步实施。

在叙事上，更稳妥的做法是把这一事件视为一条需要地方中介才能运作的链：中央的命令、地方官与精英的选择、运输与劳力的条件、以及普通家庭可能面对的迁徙、赋役或安全风险彼此相连。现有材料能够显示其中若干环节，却保留了大量沉默。承认这种沉默并非放弃解释，而是避免以一条编年记录覆盖不同区域和社群的差异。

\eventanchor{JH-LEDGER-001-6}{太康二年（281年）·“汲郡魏襄王墓出土竹书，朝廷命荀勖、和峤等整理其书写材料”的异源材料与影响范围的复核}账本条目\texttt{JH-LEDGER-001-6}记载，太康二年（281年），西晋发生““汲郡魏襄王墓出土竹书，朝廷命荀勖、和峤等整理其书写材料”的异源材料与影响范围的复核”。这一节点可放在核心节点发生后仍须处理其遗留问题的条件下理解；其后续影响被账本概括为后续影响须在不同地区和材料类型中分别核验。就地方尺度而言，事件并不只属于朝廷或统帅：征发、道路、仓储、户籍、市场、寺院与家庭关系都可能决定命令如何抵达、战事如何延长，或接收何以在不同地区呈现不一样的速度。

本条列举的核心材料为《晋书·武帝纪》；《晋书·束皙传》；《荀勖传》；汲郡魏襄王墓出土竹书。这些材料较适合钉住事件的发生、官方表述或特定人物、地点，却不自动构成对全境人口、收入、兵数、信仰或动机的调查。账本的置信与材料提示是“核心事项已有既有账本来源；本分解节点的参与者、执行范围和时间先后仍须逐案核验”；来源限度进一步说明“西晋秘阁整理、汲冢书传本与战国史书重构研究”。因此，不能把条目中的政权名称、行政分类或战报修辞直接转译为固定族群和统一社会意志，也不能将制度宣布写成各地同步实施。

在叙事上，更稳妥的做法是把这一事件视为一条需要地方中介才能运作的链：中央的命令、地方官与精英的选择、运输与劳力的条件、以及普通家庭可能面对的迁徙、赋役或安全风险彼此相连。现有材料能够显示其中若干环节，却保留了大量沉默。承认这种沉默并非放弃解释，而是避免以一条编年记录覆盖不同区域和社群的差异。

\eventanchor{JH-LEDGER-002-1}{太熙元年（289年）·“武帝病重时以外戚杨骏受遗诏辅政，围绕太子司马衷的继承安排集中于宫廷”的前期动员与资源准备}账本条目\texttt{JH-LEDGER-002-1}记载，太熙元年（289年），西晋发生““武帝病重时以外戚杨骏受遗诏辅政，围绕太子司马衷的继承安排集中于宫廷”的前期动员与资源准备”。这一节点可放在当时的权力、资源与信息约束推动相关过程展开的条件下理解；其后续影响被账本概括为这一过程为“武帝病重时以外戚杨骏受遗诏辅政，围绕太子司马衷的继承安排集中于宫廷”提供可观察的条件。就地方尺度而言，事件并不只属于朝廷或统帅：征发、道路、仓储、户籍、市场、寺院与家庭关系都可能决定命令如何抵达、战事如何延长，或接收何以在不同地区呈现不一样的速度。

本条列举的核心材料为《晋书·武帝纪》；《晋书·杨骏传》；《资治通鉴·晋纪》。这些材料较适合钉住事件的发生、官方表述或特定人物、地点，却不自动构成对全境人口、收入、兵数、信仰或动机的调查。账本的置信与材料提示是“核心事项已有既有账本来源；本分解节点的参与者、执行范围和时间先后仍须逐案核验”；来源限度进一步说明“西晋外戚辅政、遗诏政治与宫城军权研究”。因此，不能把条目中的政权名称、行政分类或战报修辞直接转译为固定族群和统一社会意志，也不能将制度宣布写成各地同步实施。

在叙事上，更稳妥的做法是把这一事件视为一条需要地方中介才能运作的链：中央的命令、地方官与精英的选择、运输与劳力的条件、以及普通家庭可能面对的迁徙、赋役或安全风险彼此相连。现有材料能够显示其中若干环节，却保留了大量沉默。承认这种沉默并非放弃解释，而是避免以一条编年记录覆盖不同区域和社群的差异。

\subsection{信仰、知识与物质空间}

\eventanchor{JH-LEDGER-002-2}{太熙元年（289年）·“武帝病重时以外戚杨骏受遗诏辅政，围绕太子司马衷的继承安排集中于宫廷”的命令、联盟或地方响应}账本条目\texttt{JH-LEDGER-002-2}记载，太熙元年（289年），西晋发生““武帝病重时以外戚杨骏受遗诏辅政，围绕太子司马衷的继承安排集中于宫廷”的命令、联盟或地方响应”。这一节点可放在当时的权力、资源与信息约束推动相关过程展开的条件下理解；其后续影响被账本概括为这一过程为“武帝病重时以外戚杨骏受遗诏辅政，围绕太子司马衷的继承安排集中于宫廷”提供可观察的条件。就地方尺度而言，事件并不只属于朝廷或统帅：征发、道路、仓储、户籍、市场、寺院与家庭关系都可能决定命令如何抵达、战事如何延长，或接收何以在不同地区呈现不一样的速度。

本条列举的核心材料为《晋书·武帝纪》；《晋书·杨骏传》；《资治通鉴·晋纪》。这些材料较适合钉住事件的发生、官方表述或特定人物、地点，却不自动构成对全境人口、收入、兵数、信仰或动机的调查。账本的置信与材料提示是“核心事项已有既有账本来源；本分解节点的参与者、执行范围和时间先后仍须逐案核验”；来源限度进一步说明“西晋外戚辅政、遗诏政治与宫城军权研究”。因此，不能把条目中的政权名称、行政分类或战报修辞直接转译为固定族群和统一社会意志，也不能将制度宣布写成各地同步实施。

在叙事上，更稳妥的做法是把这一事件视为一条需要地方中介才能运作的链：中央的命令、地方官与精英的选择、运输与劳力的条件、以及普通家庭可能面对的迁徙、赋役或安全风险彼此相连。现有材料能够显示其中若干环节，却保留了大量沉默。承认这种沉默并非放弃解释，而是避免以一条编年记录覆盖不同区域和社群的差异。

\eventanchor{JH-LEDGER-002-3}{太熙元年（289年）·武帝病重时以外戚杨骏受遗诏辅政，围绕太子司马衷的继承安排集中于宫廷}账本条目\texttt{JH-LEDGER-002-3}记载，太熙元年（289年），西晋发生“武帝病重时以外戚杨骏受遗诏辅政，围绕太子司马衷的继承安排集中于宫廷”。这一节点可放在“武帝病重时以外戚杨骏受遗诏辅政，围绕太子司马衷的继承安排集中于宫廷”发生的政治与区域条件共同作用的条件下理解；其后续影响被账本概括为该节点改变后续资源、联盟或制度处置的条件。就地方尺度而言，事件并不只属于朝廷或统帅：征发、道路、仓储、户籍、市场、寺院与家庭关系都可能决定命令如何抵达、战事如何延长，或接收何以在不同地区呈现不一样的速度。

本条列举的核心材料为《晋书·武帝纪》；《晋书·杨骏传》；《资治通鉴·晋纪》。这些材料较适合钉住事件的发生、官方表述或特定人物、地点，却不自动构成对全境人口、收入、兵数、信仰或动机的调查。账本的置信与材料提示是“核心事项已有既有账本来源；本分解节点的参与者、执行范围和时间先后仍须逐案核验”；来源限度进一步说明“西晋外戚辅政、遗诏政治与宫城军权研究”。因此，不能把条目中的政权名称、行政分类或战报修辞直接转译为固定族群和统一社会意志，也不能将制度宣布写成各地同步实施。

在叙事上，更稳妥的做法是把这一事件视为一条需要地方中介才能运作的链：中央的命令、地方官与精英的选择、运输与劳力的条件、以及普通家庭可能面对的迁徙、赋役或安全风险彼此相连。现有材料能够显示其中若干环节，却保留了大量沉默。承认这种沉默并非放弃解释，而是避免以一条编年记录覆盖不同区域和社群的差异。

\eventanchor{JH-LEDGER-002-4}{太熙元年（289年）·“武帝病重时以外戚杨骏受遗诏辅政，围绕太子司马衷的继承安排集中于宫廷”的执行中的空间与人员处置}账本条目\texttt{JH-LEDGER-002-4}记载，太熙元年（289年），西晋发生““武帝病重时以外戚杨骏受遗诏辅政，围绕太子司马衷的继承安排集中于宫廷”的执行中的空间与人员处置”。这一节点可放在“武帝病重时以外戚杨骏受遗诏辅政，围绕太子司马衷的继承安排集中于宫廷”发生的政治与区域条件共同作用的条件下理解；其后续影响被账本概括为该节点改变后续资源、联盟或制度处置的条件。就地方尺度而言，事件并不只属于朝廷或统帅：征发、道路、仓储、户籍、市场、寺院与家庭关系都可能决定命令如何抵达、战事如何延长，或接收何以在不同地区呈现不一样的速度。

本条列举的核心材料为《晋书·武帝纪》；《晋书·杨骏传》；《资治通鉴·晋纪》。这些材料较适合钉住事件的发生、官方表述或特定人物、地点，却不自动构成对全境人口、收入、兵数、信仰或动机的调查。账本的置信与材料提示是“核心事项已有既有账本来源；本分解节点的参与者、执行范围和时间先后仍须逐案核验”；来源限度进一步说明“西晋外戚辅政、遗诏政治与宫城军权研究”。因此，不能把条目中的政权名称、行政分类或战报修辞直接转译为固定族群和统一社会意志，也不能将制度宣布写成各地同步实施。

在叙事上，更稳妥的做法是把这一事件视为一条需要地方中介才能运作的链：中央的命令、地方官与精英的选择、运输与劳力的条件、以及普通家庭可能面对的迁徙、赋役或安全风险彼此相连。现有材料能够显示其中若干环节，却保留了大量沉默。承认这种沉默并非放弃解释，而是避免以一条编年记录覆盖不同区域和社群的差异。

\eventanchor{JH-LEDGER-002-5}{太熙元年（289年）·“武帝病重时以外戚杨骏受遗诏辅政，围绕太子司马衷的继承安排集中于宫廷”的后续制度或军政调整}账本条目\texttt{JH-LEDGER-002-5}记载，太熙元年（289年），西晋发生““武帝病重时以外戚杨骏受遗诏辅政，围绕太子司马衷的继承安排集中于宫廷”的后续制度或军政调整”。这一节点可放在核心节点发生后仍须处理其遗留问题的条件下理解；其后续影响被账本概括为后续影响须在不同地区和材料类型中分别核验。就地方尺度而言，事件并不只属于朝廷或统帅：征发、道路、仓储、户籍、市场、寺院与家庭关系都可能决定命令如何抵达、战事如何延长，或接收何以在不同地区呈现不一样的速度。

本条列举的核心材料为《晋书·武帝纪》；《晋书·杨骏传》；《资治通鉴·晋纪》。这些材料较适合钉住事件的发生、官方表述或特定人物、地点，却不自动构成对全境人口、收入、兵数、信仰或动机的调查。账本的置信与材料提示是“核心事项已有既有账本来源；本分解节点的参与者、执行范围和时间先后仍须逐案核验”；来源限度进一步说明“西晋外戚辅政、遗诏政治与宫城军权研究”。因此，不能把条目中的政权名称、行政分类或战报修辞直接转译为固定族群和统一社会意志，也不能将制度宣布写成各地同步实施。

在叙事上，更稳妥的做法是把这一事件视为一条需要地方中介才能运作的链：中央的命令、地方官与精英的选择、运输与劳力的条件、以及普通家庭可能面对的迁徙、赋役或安全风险彼此相连。现有材料能够显示其中若干环节，却保留了大量沉默。承认这种沉默并非放弃解释，而是避免以一条编年记录覆盖不同区域和社群的差异。

\eventanchor{JH-LEDGER-002-6}{太熙元年（289年）·“武帝病重时以外戚杨骏受遗诏辅政，围绕太子司马衷的继承安排集中于宫廷”的异源材料与影响范围的复核}账本条目\texttt{JH-LEDGER-002-6}记载，太熙元年（289年），西晋发生““武帝病重时以外戚杨骏受遗诏辅政，围绕太子司马衷的继承安排集中于宫廷”的异源材料与影响范围的复核”。这一节点可放在核心节点发生后仍须处理其遗留问题的条件下理解；其后续影响被账本概括为后续影响须在不同地区和材料类型中分别核验。就地方尺度而言，事件并不只属于朝廷或统帅：征发、道路、仓储、户籍、市场、寺院与家庭关系都可能决定命令如何抵达、战事如何延长，或接收何以在不同地区呈现不一样的速度。

本条列举的核心材料为《晋书·武帝纪》；《晋书·杨骏传》；《资治通鉴·晋纪》。这些材料较适合钉住事件的发生、官方表述或特定人物、地点，却不自动构成对全境人口、收入、兵数、信仰或动机的调查。账本的置信与材料提示是“核心事项已有既有账本来源；本分解节点的参与者、执行范围和时间先后仍须逐案核验”；来源限度进一步说明“西晋外戚辅政、遗诏政治与宫城军权研究”。因此，不能把条目中的政权名称、行政分类或战报修辞直接转译为固定族群和统一社会意志，也不能将制度宣布写成各地同步实施。

在叙事上，更稳妥的做法是把这一事件视为一条需要地方中介才能运作的链：中央的命令、地方官与精英的选择、运输与劳力的条件、以及普通家庭可能面对的迁徙、赋役或安全风险彼此相连。现有材料能够显示其中若干环节，却保留了大量沉默。承认这种沉默并非放弃解释，而是避免以一条编年记录覆盖不同区域和社群的差异。

\eventanchor{JH-LEDGER-003-1}{太熙元年四月（290年）·“司马炎去世，惠帝司马衷即位，杨骏以太傅录尚书事主持朝政”的前期动员与资源准备}账本条目\texttt{JH-LEDGER-003-1}记载，太熙元年四月（290年），西晋发生““司马炎去世，惠帝司马衷即位，杨骏以太傅录尚书事主持朝政”的前期动员与资源准备”。这一节点可放在当时的权力、资源与信息约束推动相关过程展开的条件下理解；其后续影响被账本概括为这一过程为“司马炎去世，惠帝司马衷即位，杨骏以太傅录尚书事主持朝政”提供可观察的条件。就地方尺度而言，事件并不只属于朝廷或统帅：征发、道路、仓储、户籍、市场、寺院与家庭关系都可能决定命令如何抵达、战事如何延长，或接收何以在不同地区呈现不一样的速度。

本条列举的核心材料为《晋书·武帝纪》；《晋书·惠帝纪》；《晋书·杨骏传》。这些材料较适合钉住事件的发生、官方表述或特定人物、地点，却不自动构成对全境人口、收入、兵数、信仰或动机的调查。账本的置信与材料提示是“核心事项已有既有账本来源；本分解节点的参与者、执行范围和时间先后仍须逐案核验”；来源限度进一步说明“西晋统一后继承、外戚与官僚权力转换研究”。因此，不能把条目中的政权名称、行政分类或战报修辞直接转译为固定族群和统一社会意志，也不能将制度宣布写成各地同步实施。

在叙事上，更稳妥的做法是把这一事件视为一条需要地方中介才能运作的链：中央的命令、地方官与精英的选择、运输与劳力的条件、以及普通家庭可能面对的迁徙、赋役或安全风险彼此相连。现有材料能够显示其中若干环节，却保留了大量沉默。承认这种沉默并非放弃解释，而是避免以一条编年记录覆盖不同区域和社群的差异。

\eventanchor{JH-LEDGER-003-2}{太熙元年四月（290年）·“司马炎去世，惠帝司马衷即位，杨骏以太傅录尚书事主持朝政”的命令、联盟或地方响应}账本条目\texttt{JH-LEDGER-003-2}记载，太熙元年四月（290年），西晋发生““司马炎去世，惠帝司马衷即位，杨骏以太傅录尚书事主持朝政”的命令、联盟或地方响应”。这一节点可放在当时的权力、资源与信息约束推动相关过程展开的条件下理解；其后续影响被账本概括为这一过程为“司马炎去世，惠帝司马衷即位，杨骏以太傅录尚书事主持朝政”提供可观察的条件。就地方尺度而言，事件并不只属于朝廷或统帅：征发、道路、仓储、户籍、市场、寺院与家庭关系都可能决定命令如何抵达、战事如何延长，或接收何以在不同地区呈现不一样的速度。

本条列举的核心材料为《晋书·武帝纪》；《晋书·惠帝纪》；《晋书·杨骏传》。这些材料较适合钉住事件的发生、官方表述或特定人物、地点，却不自动构成对全境人口、收入、兵数、信仰或动机的调查。账本的置信与材料提示是“核心事项已有既有账本来源；本分解节点的参与者、执行范围和时间先后仍须逐案核验”；来源限度进一步说明“西晋统一后继承、外戚与官僚权力转换研究”。因此，不能把条目中的政权名称、行政分类或战报修辞直接转译为固定族群和统一社会意志，也不能将制度宣布写成各地同步实施。

在叙事上，更稳妥的做法是把这一事件视为一条需要地方中介才能运作的链：中央的命令、地方官与精英的选择、运输与劳力的条件、以及普通家庭可能面对的迁徙、赋役或安全风险彼此相连。现有材料能够显示其中若干环节，却保留了大量沉默。承认这种沉默并非放弃解释，而是避免以一条编年记录覆盖不同区域和社群的差异。

\eventanchor{JH-LEDGER-003-3}{太熙元年四月（290年）·司马炎去世，惠帝司马衷即位，杨骏以太傅录尚书事主持朝政}账本条目\texttt{JH-LEDGER-003-3}记载，太熙元年四月（290年），西晋发生“司马炎去世，惠帝司马衷即位，杨骏以太傅录尚书事主持朝政”。这一节点可放在“司马炎去世，惠帝司马衷即位，杨骏以太傅录尚书事主持朝政”发生的政治与区域条件共同作用的条件下理解；其后续影响被账本概括为该节点改变后续资源、联盟或制度处置的条件。就地方尺度而言，事件并不只属于朝廷或统帅：征发、道路、仓储、户籍、市场、寺院与家庭关系都可能决定命令如何抵达、战事如何延长，或接收何以在不同地区呈现不一样的速度。

本条列举的核心材料为《晋书·武帝纪》；《晋书·惠帝纪》；《晋书·杨骏传》。这些材料较适合钉住事件的发生、官方表述或特定人物、地点，却不自动构成对全境人口、收入、兵数、信仰或动机的调查。账本的置信与材料提示是“核心事项已有既有账本来源；本分解节点的参与者、执行范围和时间先后仍须逐案核验”；来源限度进一步说明“西晋统一后继承、外戚与官僚权力转换研究”。因此，不能把条目中的政权名称、行政分类或战报修辞直接转译为固定族群和统一社会意志，也不能将制度宣布写成各地同步实施。

在叙事上，更稳妥的做法是把这一事件视为一条需要地方中介才能运作的链：中央的命令、地方官与精英的选择、运输与劳力的条件、以及普通家庭可能面对的迁徙、赋役或安全风险彼此相连。现有材料能够显示其中若干环节，却保留了大量沉默。承认这种沉默并非放弃解释，而是避免以一条编年记录覆盖不同区域和社群的差异。

\eventanchor{JH-LEDGER-003-4}{太熙元年四月（290年）·“司马炎去世，惠帝司马衷即位，杨骏以太傅录尚书事主持朝政”的执行中的空间与人员处置}账本条目\texttt{JH-LEDGER-003-4}记载，太熙元年四月（290年），西晋发生““司马炎去世，惠帝司马衷即位，杨骏以太傅录尚书事主持朝政”的执行中的空间与人员处置”。这一节点可放在“司马炎去世，惠帝司马衷即位，杨骏以太傅录尚书事主持朝政”发生的政治与区域条件共同作用的条件下理解；其后续影响被账本概括为该节点改变后续资源、联盟或制度处置的条件。就地方尺度而言，事件并不只属于朝廷或统帅：征发、道路、仓储、户籍、市场、寺院与家庭关系都可能决定命令如何抵达、战事如何延长，或接收何以在不同地区呈现不一样的速度。

本条列举的核心材料为《晋书·武帝纪》；《晋书·惠帝纪》；《晋书·杨骏传》。这些材料较适合钉住事件的发生、官方表述或特定人物、地点，却不自动构成对全境人口、收入、兵数、信仰或动机的调查。账本的置信与材料提示是“核心事项已有既有账本来源；本分解节点的参与者、执行范围和时间先后仍须逐案核验”；来源限度进一步说明“西晋统一后继承、外戚与官僚权力转换研究”。因此，不能把条目中的政权名称、行政分类或战报修辞直接转译为固定族群和统一社会意志，也不能将制度宣布写成各地同步实施。

在叙事上，更稳妥的做法是把这一事件视为一条需要地方中介才能运作的链：中央的命令、地方官与精英的选择、运输与劳力的条件、以及普通家庭可能面对的迁徙、赋役或安全风险彼此相连。现有材料能够显示其中若干环节，却保留了大量沉默。承认这种沉默并非放弃解释，而是避免以一条编年记录覆盖不同区域和社群的差异。

\eventanchor{JH-LEDGER-003-5}{太熙元年四月（290年）·“司马炎去世，惠帝司马衷即位，杨骏以太傅录尚书事主持朝政”的后续制度或军政调整}账本条目\texttt{JH-LEDGER-003-5}记载，太熙元年四月（290年），西晋发生““司马炎去世，惠帝司马衷即位，杨骏以太傅录尚书事主持朝政”的后续制度或军政调整”。这一节点可放在核心节点发生后仍须处理其遗留问题的条件下理解；其后续影响被账本概括为后续影响须在不同地区和材料类型中分别核验。就地方尺度而言，事件并不只属于朝廷或统帅：征发、道路、仓储、户籍、市场、寺院与家庭关系都可能决定命令如何抵达、战事如何延长，或接收何以在不同地区呈现不一样的速度。

本条列举的核心材料为《晋书·武帝纪》；《晋书·惠帝纪》；《晋书·杨骏传》。这些材料较适合钉住事件的发生、官方表述或特定人物、地点，却不自动构成对全境人口、收入、兵数、信仰或动机的调查。账本的置信与材料提示是“核心事项已有既有账本来源；本分解节点的参与者、执行范围和时间先后仍须逐案核验”；来源限度进一步说明“西晋统一后继承、外戚与官僚权力转换研究”。因此，不能把条目中的政权名称、行政分类或战报修辞直接转译为固定族群和统一社会意志，也不能将制度宣布写成各地同步实施。

在叙事上，更稳妥的做法是把这一事件视为一条需要地方中介才能运作的链：中央的命令、地方官与精英的选择、运输与劳力的条件、以及普通家庭可能面对的迁徙、赋役或安全风险彼此相连。现有材料能够显示其中若干环节，却保留了大量沉默。承认这种沉默并非放弃解释，而是避免以一条编年记录覆盖不同区域和社群的差异。

\eventanchor{JH-LEDGER-003-6}{太熙元年四月（290年）·“司马炎去世，惠帝司马衷即位，杨骏以太傅录尚书事主持朝政”的异源材料与影响范围的复核}账本条目\texttt{JH-LEDGER-003-6}记载，太熙元年四月（290年），西晋发生““司马炎去世，惠帝司马衷即位，杨骏以太傅录尚书事主持朝政”的异源材料与影响范围的复核”。这一节点可放在核心节点发生后仍须处理其遗留问题的条件下理解；其后续影响被账本概括为后续影响须在不同地区和材料类型中分别核验。就地方尺度而言，事件并不只属于朝廷或统帅：征发、道路、仓储、户籍、市场、寺院与家庭关系都可能决定命令如何抵达、战事如何延长，或接收何以在不同地区呈现不一样的速度。

本条列举的核心材料为《晋书·武帝纪》；《晋书·惠帝纪》；《晋书·杨骏传》。这些材料较适合钉住事件的发生、官方表述或特定人物、地点，却不自动构成对全境人口、收入、兵数、信仰或动机的调查。账本的置信与材料提示是“核心事项已有既有账本来源；本分解节点的参与者、执行范围和时间先后仍须逐案核验”；来源限度进一步说明“西晋统一后继承、外戚与官僚权力转换研究”。因此，不能把条目中的政权名称、行政分类或战报修辞直接转译为固定族群和统一社会意志，也不能将制度宣布写成各地同步实施。

在叙事上，更稳妥的做法是把这一事件视为一条需要地方中介才能运作的链：中央的命令、地方官与精英的选择、运输与劳力的条件、以及普通家庭可能面对的迁徙、赋役或安全风险彼此相连。现有材料能够显示其中若干环节，却保留了大量沉默。承认这种沉默并非放弃解释，而是避免以一条编年记录覆盖不同区域和社群的差异。
